\chapter{Bones and Muscles}\label{ch:g3:bones}

Stand up. Bend your \omterm{def:g1:my-body:parts}{arm}. Take a step. Under your \omterm{def:g1:senses:five}{skin}, \omterm{def:g3:bones:bone}{bones} and
\omterm{def:g3:bones:muscle}{muscles} work as a team. \omterm{def:g3:bones:bone}{Bones} give shape and support; \omterm{def:g3:bones:muscle}{muscles} pull
to make movement.

\section{The skeleton}

\begin{definition}[Bone and skeleton]\label{def:g3:bones:bone}
A \emph{bone}\index{bone} is a hard organ that supports the \omterm{def:g1:my-body:body}{body}.
All the bones together form the \emph{skeleton}\index{skeleton}.
The skeleton holds you up, protects soft organs, and gives \omterm{def:g3:bones:muscle}{muscles}
\omterm{def:g1:caring:needs}{places} to pull.
\end{definition}

\begin{omfigure}
\begin{tikzpicture}[font=\footnotesize]
  % stick skeleton
  \draw[thick, omDef] (0,4.0) circle [radius=0.4]; % skull
  \draw[thick, omDef] (0,3.55) -- (0,2.0); % spine
  \draw[thick, omThm] (-1.1,3.1) -- (0,2.9) -- (1.1,3.1); % arms
  \draw[thick, omThm] (-1.1,3.1) -- (-1.5,2.4);
  \draw[thick, omThm] (1.1,3.1) -- (1.5,2.4);
  \draw[thick, omProp] (-0.5,0.2) -- (0,2.0) -- (0.5,0.2); % legs
  \draw[thick, omProp] (-0.5,0.2) -- (-0.7,-0.3);
  \draw[thick, omProp] (0.5,0.2) -- (0.7,-0.3);
  \node[omDef, right] at (0.5,4.0) {skull};
  \node[omDef, right] at (0.15,2.7) {spine};
  \node[omThm, left] at (-1.2,3.2) {arm bones};
  \node[omProp, right] at (0.6,0.8) {leg bones};
\end{tikzpicture}

{\small Simplified \omterm{def:g3:bones:bone}{skeleton} map. Real \omterm{def:g3:bones:bone}{bones} are many and shaped for
their jobs --- more than 200 in an \omterm{def:g3:life-cycles:butterfly}{adult} human.}
\end{omfigure}

\begin{example}[Protection]\label{ex:g3:bones:protect}
The skull protects the brain. Ribs protect the heart and lungs.
The backbone protects the spinal cord and holds the \omterm{def:g1:my-body:body}{body} upright.
\end{example}

\begin{definition}[Joint]\label{def:g3:bones:joint}
A \emph{joint}\index{joint} is where two \omterm{def:g3:bones:bone}{bones} meet. Joints let the
\omterm{def:g3:bones:bone}{skeleton} bend --- knees, elbows, fingers, neck.
\end{definition}

\section{Muscles}

\begin{definition}[Muscle]\label{def:g3:bones:muscle}
A \emph{muscle}\index{muscle} is a soft tissue that can contract
(get shorter) to pull on \omterm{def:g3:bones:bone}{bones}. Muscles cannot push --- they only
pull. That is why they often work in pairs.
\end{definition}

\begin{omfigure}
\begin{tikzpicture}[font=\footnotesize]
  % arm bent
  \draw[very thick, omDef] (0,0) -- (0,2.0) -- (1.5,2.8);
  \fill[omThm!40] (-0.35,0.8) ellipse (0.35 and 0.7);
  \node[omThm, left] at (-0.8,0.9) {muscle};
  \draw[-Stealth, thick, omExo] (0.3,1.5) -- (0.9,2.2);
  \node[omExo, right] at (0.5,1.3) {pulls};
  \node at (0.5,-0.5) {muscle pulls bone $\to$ arm bends};
\end{tikzpicture}

{\small \omterm{def:g3:bones:muscle}{Muscles} pull on \omterm{def:g3:bones:bone}{bones} across a \omterm{def:g3:bones:joint}{joint} to bend or straighten
a limb.}
\end{omfigure}

\begin{method}[Feel them work]\label{met:g3:bones:feel}
\begin{enumerate}
  \item Put your left \omterm{def:g1:my-body:parts}{hand} on your right upper \omterm{def:g1:my-body:parts}{arm}.
  \item Bend your right elbow: feel the front \omterm{def:g3:bones:muscle}{muscle} firm up.
  \item Straighten the \omterm{def:g1:my-body:parts}{arm}: a \omterm{def:g3:bones:muscle}{muscle} on the back works.
  \item \omterm{def:g3:bones:bone}{Bones} are the levers; \omterm{def:g3:bones:muscle}{muscles} are the motors.
\end{enumerate}
\end{method}

\begin{example}[Pairs]\label{ex:g3:bones:pairs}
To bend the elbow, one \omterm{def:g3:bones:muscle}{muscle} shortens. To straighten it, the
partner \omterm{def:g3:bones:muscle}{muscle} shortens while the first relaxes. Teamwork!
\end{example}

\section{Exercises}

\begin{exercise}[$\star$]\label{exo:g3:bones:1}
What is the \omterm{def:g3:bones:bone}{skeleton}?
\end{exercise}

\begin{exercise}[$\star$]\label{exo:g3:bones:2}
Name two jobs of \omterm{def:g3:bones:bone}{bones}.
\end{exercise}

\begin{exercise}[$\star$]\label{exo:g3:bones:3}
What is a \omterm{def:g3:bones:joint}{joint}?
\end{exercise}

\begin{exercise}[$\star$]\label{exo:g3:bones:4}
True or false? ``\omterm{def:g3:bones:muscle}{Muscles} can push \omterm{def:g3:bones:bone}{bones}.''
\end{exercise}

\begin{exercise}[$\star\star$]\label{exo:g3:bones:5}
Why do \omterm{def:g1:my-body:parts}{arm} \omterm{def:g3:bones:muscle}{muscles} work in pairs?
\end{exercise}

